
\section{Planetary Orbits}
 \begin{table}[h]
 \begin{tabular}{|l| l| l| l| l| l| l|}
 \hline
name & 
\multicolumn{1}{|p{1.5cm}|}{semi-major axis} & 
      eccentricity &
     \multicolumn{1}{|p{1.5cm}|}{inclination to ecliptic} &
     \multicolumn{1}{|p{1.5cm}|}{longitude of \mbox{ascending} node} &
     \multicolumn{1}{|p{2.0cm}|}{argument of perihelion} &
     \multicolumn{1}{|p{2.cm}|}{time of \mbox{perihelion}}\\
\hline
Mercury & 0.3487 AU & 0.206 & 7.005°  & 48.331° & 29.124° & 2024-Dec-25\\
Venus&&&&&&\\
Earth&149598023 km & 0.0167086& ---   &−11.26064°&114.20783°&2023-Jan-04\\
Mars    & 1.524 AU & 0.0934 & 1.850° &49.57854° & 286.5°   &June 21, 2022\\
Jupiter & 5.204 AU & 0.0489 & 1.303° & 100.464° & 273.867° & January 20, 2023\\
Saturn&&&&&&\\
Uranus&&&&&&\\
Neptune&&&&&&\\
Pluto&&&&&&\\
\hline
\end{tabular}
\begin{tabular}{|l| l| l| l| l| l| l|}
\hline
name & 
\multicolumn{1}{|p{1.5cm}|}{semi-major axis} & 
eccentricity &
\multicolumn{1}{|p{1.5cm}|}{inclination to ecliptic} &
\multicolumn{1}{|p{1.5cm}|}{longitude of \mbox{ascending} node} &
\multicolumn{1}{|p{2.0cm}|}{argument of perihelion} &
\multicolumn{1}{|p{2.cm}|}{time of \mbox{perihelion}}\\
\hline
Mercury & 0.3871 AU & 0.2056 & 7.005°  & 48.331° & 29.124° & 2024-Dec-25\\
Venus   & 0.7233 AU & 0.0068 & 3.394°  & 76.680° & 54.884° & 2023-Dec-31\\
Earth   & 149598023 km & 0.0167086& ---   &−11.26064°&114.20783°&2023-Jan-04\\
Mars    & 1.5237 AU & 0.0934 & 1.850° &49.57854° & 286.5°   &2022-Jun-21\\
Jupiter & 5.204 AU & 0.0489 & 1.303° & 100.464° & 273.867° & 2023-Jan-20\\
Saturn  & 9.582 AU & 0.0565 & 2.485° & 113.665° & 339.392° & 2003-Jul-26\\
Uranus  & 19.201 AU & 0.0457 & 0.773° & 74.006° & 96.998° & 1966-Apr-14\\
Neptune & 30.047 AU & 0.0086 & 1.770° & 131.784° & 273.187° & 2042-Sep-05\\
Pluto   & 39.482 AU & 0.2488 & 17.140° & 110.299° & 113.834° & 1989-Sep-05\\
\hline
\end{tabular}

\caption{Tables of planetray orbits}
\end{table}




\section{Attributes of Jupiter}
\subsection{Data Model}


\begin{table}
\begin{tabular}{ | l | l |}
\hline
\multicolumn{2}{|l|}{Physical Characteristics}  \\
\hline
Equatorial radius&71492 km  \\
Polar radius&66854 km  \\
Mass&1.8982×$10^27$ kg  \\
Sidereal rotation period&9.9250 hours  \\
North pole right ascension&268.057° \\
North pole declination&64.495° \\
Absolute magnitude (H)&−9.4 \\
Moment of inertia factor&0.2756±0.0006  \\
\hline
\end{tabular}

\caption{Chosen data model for physical characteristics}
\end{table}

\newpage
\subsection{Workings}
Start with text\\

\begin{minipage}{\textwidth}  %minipage so as to catch footnotes
\begin{tabular}{ | l | l | p{3cm}|}
\hline
\multicolumn{3}{|l|}{Orbital Characteristics}  \\
\hline
Aphelion&5.4570 AU (816.363 million km)        & \multirow{4}{3cm}{any two of these 
                                                                   four to specify the size and shape of the elipse}\\
Perihelion&4.9506 AU (740.595 million km)      &  \\
Semi-major axis\footnote{could 
instead use either of orbital period(synodic or sideral) or average orbital speed} 
                                               &5.2038 AU (778.479 million km) &  \\
Eccentricity&0.0489                            &  \\
\hline
Inclination to ecliptic\footnote{could instead use `inclination to sun\textquotesingle s equator' or `inclination to invariable plane'}& 1.303° & \multirow{2}{3cm}{these two specify the 
                                                               position of the orbital plane in space 
                                                                }\\
Longitude of ascending node&100.464° & \\[1cm]
\hline
Argument of perihelion&273.867° & This specifies the position of the elipse within the orbital plane\\
\hline
Time of perihelion\footnote{could instead use `mean anomaly' on a reference time} &January 21, 2023 & this positions the planet within its orbit \\
\hline
\end{tabular}
\end{minipage}

\begin{tabular}{ | l | l | p{3cm}|}
\hline
Physical characteristics & &  \\
\hline
Mean radius&69911 km & calculated\\
Equatorial radius&71492 km & \\
Polar radius&66854 km & \\
Flattening&0.06487 & calculated\\
Surface area&6.1469×1010 km2 & calculated\\
Volume&1.4313×1015 km3 & calculated\\
Mass&1.8982×1027 kg & \\
Mean density&1.326 g/cm3 & calculated\\
Surface gravity&24.79 m/s2 & calculated\\
Moment of inertia factor&0.2756±0.0006 & \\
Escape velocity&59.5 km/s & calculated\\
Synodic rotation period&9.9258 h & calculate from sideral rotation period \\
Sidereal rotation period&9.9250 hours  & \\
Equatorial rotation velocity&12.6 km/s & calculated\\
Axial tilt&3.13° (to orbit) & calculate this\\
North pole right ascension&268.057°; 17h 52m 14s & \\
North pole declination&64.495° & \\
Albedo&&\\  
&0.503 (Bond)  & \\
&0.538 (geometric) & can be calculated from diameter and absolute magnitude\\
Temperature&88 K (−185 °C)  & (blackbody temperature) \\
Surface temp.&min mean max  &\\
1 bar& 165 K & \\
0.1 bar&78 K, 128 K& \\
Apparent magnitude&−2.94[16] to −1.66 & \\
Absolute magnitude (H)&−9.4 & \\
Angular diameter&29.8" to 50.1" & \\
\hline
\end{tabular}

and so on and on and on